\section{Questions raised at the sitting}
\label{sec:options}

Each subsection states a question raised by a coauthor, what Gate~0
currently does, what the alternative would require, and the measurement
that would distinguish them. No option is selected here.

It is worth separating two classes of alternative. Some are
reparameterisations or single-equation substitutions within the two-week
recursion of \S\ref{sec:dynamics}. Others replace the host: per-agent
constrained optimisation as in Agrimate \citep{kuhla2025}, or annual
spatial price equilibrium \citep{samuelson1952,takayama1971}. The second
class is not an edit. Figure~\ref{fig:pways} shows that division.

\subsection{Which object is ``the baseline''?}

The term was used for four distinct objects, and the ambiguity is worth
resolving before anything else.

The \emph{reference price} \(p_0\) is the start-year mean of the World
Bank Pink Sheet in constant 2010 dollars, where \(p_0\) enters
\eqref{eq:demand} as the price at which \(C^{\mathrm{flex}}_{i,t}\) is
defined and \eqref{eq:pscar} as the level the scarcity price returns to.
The \emph{climatological harvest} \(H^{\mathrm{seas}}_{i,t}\) is
window-mean production distributed over the calendar, and enters
expectations through \eqref{eq:hexp}. The \emph{spin-up} initialises
stocks \(S_{i,t}\) from 2005 USDA levels and runs the unshocked recursion
for two years, so that opening stocks are dynamically consistent rather
than a raw data snapshot. The \emph{twin} is a full parallel run of the
same recursion under \(H^{\mathrm{seas}}\), mean food and feed demand, no
industrial excess, and \(\tau_{i,t}=0\), which supplies
\(\mathcal{F}^{\mathrm{twin}}_t\) in \eqref{eq:ratio}. Scarcity is measured
against the twin, not against \(p_0\) alone.

Agrimate spins up for roughly four years. Lengthening ours is a
sensitivity test on initial conditions, not a change of theory. Replacing
the twin with an annual market-clearing reference is a change of host
(Fig.~\ref{fig:pways}D).

\flowfig{02_baselines_four_objects.pdf}{The four objects that have all been
called ``baseline.'' The scarcity term in \eqref{eq:pscar} is measured
against the twin. \label{fig:b0}}

\subsection{Why is available surplus not fully offered each period?}

Because \eqref{eq:offers} withholds the target stock \(T_{i,t}\) defined in
\eqref{eq:target}. A country consumes \(d_{i,t}\), retains lean-season
cover plus safety stock, and offers the residual. No interest rate and no
expected-price comparison appear: there is no condition of the form
\(\mathbb{E}_t[p_{t+1}]\ge (1+r)p_t + c\) that a competitive storer would
impose, where \(r\) is the interest rate and \(c\) the physical cost of
carry \citep{wright1982,williams1991,deaton1992}. Slide~48 of the deck is
the offer-price update \eqref{eq:ask}, not a cost of carry; that reading
should be confirmed with the authors.

Agrimate's commercial supplier does solve a store-versus-sell problem from
expected profit, conditional on announced restrictions
\citep{kuhla2025}. Adopting that is Fig.~\ref{fig:pways}C, not a
modification of \eqref{eq:target}. Note also that Agrimate separates
commercial storage from a strategic-reserve agent following a
partial-adjustment refill rule; SHEAF's Gate~2 food-security trigger is not
that second agent, and the two should not be equated when comparing the
storage layers. The intermediate option is an explicit
arbitrage margin: hold above \(T_{i,t}\) only when
\(\mathbb{E}_t[p_{t+1}]-p_t\) exceeds a per-step carry cost. This
introduces both a price expectation and a carry parameter, neither of
which the model currently has. A related experiment --- country-specific
safety stock ratios \(\sigma^{\mathrm{stu}}_i\) --- improved US wheat
stocks but inflated Chinese maize stocks by roughly a factor of three, and
was not retained.

\flowfig{10_storage_cover_rule.pdf}{Partition of available grain under
\eqref{eq:offers}: consumption, cover, offer. No store-versus-sell
programme is solved. \label{fig:s0}}

\subsection{How is the world price determined?}

By \eqref{eq:pscar}--\eqref{eq:pupdate}: a shipment-weighted offer price
combined with a scarcity term and then smoothed. The system is a discrete
map, not a differential equation; the only relaxation-like element is the
offer-price adjustment \eqref{eq:ask}.

Three distinct requests were compressed into ``there should be an
optimisation.'' The first is \emph{agent optimisation}: per-region
programmes at every step, i.e.\ Agrimate (Fig.~\ref{fig:pways}C). The
second is \emph{market clearing}: price as the multiplier on a balance
constraint in an annual spatial equilibrium (Fig.~\ref{fig:pways}D) ---
the host that produced the wrong sign for 2007/08. The third is the
observation that Gate~0 may already contain an optimisation principle,
since \eqref{eq:pscar} has the form of inverse demand over accessible
stocks with elasticity \(\eta\).

Three gaps remain between that form and a genuine equilibrium, and they
should be stated rather than papered over. Demand \eqref{eq:demand} uses
\(p_{t-1}\), so quantities entering trade are not consistent with the
price the same period produces. Blending and smoothing in
\eqref{eq:pstar}--\eqref{eq:pupdate} means \(p_t\) does not lie on the
inverse-demand schedule within the period. And \eqref{eq:target} is a
target-stock rule, so storage is not a first-order condition. Writing
these three sentences accurately is cheaper, and more defensible, than
attaching an \(\arg\max\) to the existing recursion. TWIST likewise did
not solve a storage Euler equation \citep{schewe2017}.

Two alternatives remain inside the two-week host. The first is
contemporaneous consistency: replace \(p_{t-1}\) with \(p_t\) in
\eqref{eq:demand} and solve the resulting scalar fixed point
\(p_t=\mathcal{G}(p_t)\) each period, where \(\mathcal{G}\) is the
composition of \eqref{eq:demand}--\eqref{eq:pupdate}. That is one
univariate root-find per crop per step, not a system of programmes. The
second is a structural offer rule: replace the exponential adjustment in
\eqref{eq:ask} with a markup over marginal cost, retaining the
inverse-demand-plus-offers structure of \eqref{eq:pstar}.

\flowfig{20_four_ways_p.pdf}{Four hosts capable of producing a world
price. A is the current recursion. B remains the two-week host with
demand and price made contemporaneously consistent. C and D are different
models. \label{fig:pways}}

\subsection{Who trades with whom, and at which price?}

Traded quantities are determined by \eqref{eq:offers} and
\eqref{eq:importdemand}; the FAOSTAT matrices \(A_{ij}\) and \(S_{ij}\)
supply only the bilateral pattern, reweighted by \eqref{eq:reweight}. A
restriction reduces \(O_{i,t}\) directly, whereas realised shipments
\(X_{ij,t}\) are the short side of \eqref{eq:clear}; Russian wheat offers
in late 2010 therefore collapse by far more than Russian shipments do.
The phrase ``sold to each market'' in the notes most likely refers to
Agrimate's purchaser agent rather than to its storer.

A restriction does not currently generate a domestic price below \(p_t\).
Introducing a domestic price \(p^{\mathrm{dom}}_{i,t}\) distinct from the
world price would matter for a consumer-welfare analysis and would connect
to the price-insulation literature \citep{martin2012}; it is not required
to reproduce the single World Bank world price series that Gate~0 is
scored against.

\subsection{Is demand linear, and what counts as industrial use?}

Demand is isoelastic, not linear: \eqref{eq:demand} is
\(C^{\mathrm{flex}}(p_{t-1}/p_0)^{\varepsilon}\), so the elasticity
\(\varepsilon\) is constant and the price response is multiplicative. The
linear system \(D=a-Mp\) belongs to the parked annual host, where it is the
quadratic-programme demand block; it is not the crisis specification. The
two should not be conflated when comparing SHEAF against the annual
formulation.

The industrial component \(C^{\mathrm{ind}}_{i,t}\) is constructed as US
maize food-seed-industrial use minus its 2000--04 mean, distributed evenly
across periods, and is zero for every other country and crop. It is the
post-Renewable-Fuel-Standard ethanol expansion treated as an exogenous,
price-inelastic demand shifter \citep{rfs2005,abbott2011,roberts2013}, not
an endogenous biofuel sector.

For official scores, \(C^{\mathrm{flex}}_{i,t}\) is the window-mean food
and feed level divided by 24 rather than the year-by-year USDA path. That
choice matches Agrimate's treatment and is the reason simulated world
consumption does not track the USDA series; using the observed annual path
is available as a sensitivity but would change what the official split
means.

\subsection{Should a surplus lower the price, and on what frequency?}

It does, within the same period. The exponent \(\eta\) in \eqref{eq:pscar}
is symmetric: when accessible stocks exceed the twin, \(r_t<1\) and the
scarcity price falls below \(p_0\), which raises food and feed demand
through \eqref{eq:demand} in the following period. An earlier
specification muted the abundance side, and the consequence was that
surpluses accumulated in storage instead of clearing.

The frequency question is separate and worth stating plainly. Abundance in
grain markets is resolved over a marketing year, not a fortnight, and
Gate~0 has no mechanism that operates at that frequency: there is no
acreage response, and \(\phi\) in \eqref{eq:hexp} looks forward only to the
next harvest pulse. What the two-week recursion can represent is the
within-season absorption of a surplus through consumption and stock
accumulation. Multi-year adjustment is outside its scope, and neither TWIST
\citep{schewe2017} nor Agrimate \citep{kuhla2025} makes production
endogenous either.

\subsection{What is the expectational horizon?}

Short, and confined to one place. Equation \eqref{eq:hexp} places weight
\(\phi\approx 0.5\) on the current harvest realisation and the remainder on
climatology, and this expectation enters only the cover target
\eqref{eq:target}. Demand \eqref{eq:demand} uses the realised lagged price
\(p_{t-1}\), which is backward-looking rather than expectational. The
lean-season horizon in \(L_{i,t}\) is defined by arrival of 12\% of the
annual \emph{world} harvest and is therefore common across countries, not
indexed to national harvest calendars.

The statement ``prices stay near the reference unless new information
arrives'' is the identity established in \S\ref{sec:dynamics}: with
\(\mathcal{F}_t=\mathcal{F}^{\mathrm{twin}}_t\), \(\Delta u_t=0\) and
\(b_t=0\), \eqref{eq:pscar} returns \(p_0\) exactly. One live option is to
index \(L_{i,t}\) to national harvest calendars. Rational-expectations
storage over a long horizon \citep{deaton1992} is the opposite of what was
requested and would change the host.

\subsection{Is the parameter count defensible?}

The concern is legitimate and should be answered by classification and by
ablation rather than by assertion. Table~\ref{tab:params} labels each
parameter as literature-sourced (L), definitional (S), or reduced-form
(R). Reduced-form parameters have economically signed effects but are not
structural elasticities, and they are not retuned crisis by crisis.

Separately, the model contains several inequality constraints --- storage
capacity \(W_{i,t}\) in \eqref{eq:capacity}, the offer-price bounds on
\eqref{eq:ask}, the world-price bounds on \eqref{eq:pupdate}, and the AMIS
mapping to \(\tau_{i,t}\). Coauthors observed that a sufficiently
constrained rule-based system can mimic the outcome of an optimisation.
The cheapest empirical response is channel ablation: set \(\omega=1\) or
\(\omega=0\) in \eqref{eq:pstar}, set \(\alpha_r=0\) or \(\kappa_b=0\) in
\eqref{eq:ask} and \eqref{eq:pscar}, and record which channel carries
2007/08 versus 2010/11. That is a decomposition, not a refit.
